% SOURCE_AUTHORITY: sga5_fr_workpass.tex, lines 306--325 (LF-normalized unit).
% SOURCE_SCAN: SGA5 (1).pdf, Exposé I opening pages; scan page mapping pending frozen unit manifest.
\begin{center}
{\Large \textbf{SGA 5 --- Exposición I}}\\[0.4em]
{\large \textbf{COMPLEJOS DUALIZANTES}}\\[0.3em]
por A.~Grothendieck\\
(redacción de L.~Illusie)
\end{center}

\section*{Introducción}

Esta exposición está dedicada al teorema de bidualidad en cohomología étale. La teoría desarrollada aquí y en SGA~A~XVIII (SGA~A = SGA~4), para la topología étale y los haces constructibles de torsión sobre los esquemas, es formalmente muy análoga a la desarrollada en HARTSHORNE~[H] para la topología de Zariski y los haces coherentes: esta última le sirvió, en muy gran medida, de modelo. Sin embargo, mientras que en el caso de los coeficientes ``continuos'' la existencia de complejos dualizantes no presenta dificultades, en el de los coeficientes ``discretos'' resulta mucho más delicada de demostrar, y parece requerir de manera esencial la resolución de singularidades. Véase, no obstante, el teorema de bidualidad de Deligne (SGA~$4^{1/2}$~Teorema de finitud~4.3.).

El parágrafo~1 contiene la definición de los complejos dualizantes y las propiedades formales que de ella se deducen, en particular la interpretación de los funtores dualizantes como ``intercambiadores de funtores''.

En el parágrafo~2 se demuestra que, sobre un esquema localmente noetheriano y conexo, un complejo dualizante queda determinado de manera única, salvo una traslación de grados y el producto tensorial con un haz invertible.

En el parágrafo~3 se llega al teorema central de la exposición (3.4.1.), que asegura que, sobre un ``buen'' esquema regular $X$, por ejemplo un esquema excelente, noetheriano y regular de característica cero\footnote{Esta última restricción será innecesaria una vez que se disponga, para estos esquemas, de la resolución de singularidades al estilo de Hironaka y del ``teorema de pureza''.}, el complejo $(\Z/n\Z)_X$ es dualizante.

El parágrafo~4 presenta las aplicaciones de la teoría al teorema de dualidad local. Este expresa que, sobre un ``buen'' esquema $X$ provisto de un punto cerrado $x$ cuyo cuerpo residual es separablemente cerrado, y para un haz constructible $F$, se tiene un apareamiento perfecto entre $\underline{\H}^i(F')_x$ y $\underline{\H}^{-i}_x(F)$, donde $F'$ es el dual de $F$, es decir, $F' = \R\,\underline{\Hom}(F,K_X)$, y $K_X$ es un complejo dualizante.

Por último, en el parágrafo~5 se demuestra directamente que, sobre un esquema regular de dimensión~1, el complejo $(\Z/n\Z)_X$ es dualizante ($n$ primo con las características residuales).
